%% LyX 2.4.4 created this file.  For more info, see https://www.lyx.org/.
%% Do not edit unless you really know what you are doing.
\documentclass[english]{article}
\usepackage[T1]{fontenc}
\usepackage[latin9]{inputenc}
\usepackage{units}
\usepackage{enumitem}
\usepackage{amsmath}
\usepackage{amsthm}
\usepackage{amssymb}
\usepackage{cancel}
\usepackage{geometry}
\geometry{verbose,tmargin=1in,bmargin=1in,lmargin=1in,rmargin=1in}

\makeatletter
%%%%%%%%%%%%%%%%%%%%%%%%%%%%%% Textclass specific LaTeX commands.
\numberwithin{equation}{section}
\numberwithin{figure}{section}
\newlength{\lyxlabelwidth}      % auxiliary length 

%%%%%%%%%%%%%%%%%%%%%%%%%%%%%% User specified LaTeX commands.
\usepackage{fancyhdr}
\usepackage{lastpage}
\pagestyle{fancy}
\usepackage{enumitem}
\usepackage{ifthen}
%sets math fonts to be more readable
\everymath=\expandafter{\the\everymath\displaystyle}
%\DeclareMathSizes{10}{10}{10}{10}
\usepackage{multicol}

\usepackage{tikz}
\usetikzlibrary{plotmarks}
\usepackage{pgfplots}
\pgfplotsset{compat=newest}

\usepackage{calc}
\usepackage[nomessages]{fp}

\usepackage{advdate}    % Advancing/saving dates
\usepackage{datetime}   % Dates formatting
\usepackage{datenumber} % Counters for dates
\usdate
% Added by lyx2lyx
\setlength{\parskip}{\smallskipamount}
\setlength{\parindent}{0pt}

\makeatother

\usepackage{babel}
\begin{document}
\renewcommand{\author}{Dr.\,James Ashe}

\newcommand{\classNum}{131}

\newcommand{\classSec}{B}

\newcommand{\nameType}{Solutions}

\newcommand{\examType}{Exam 2: 1.7, 2.1-2.4}

\renewcommand{\date}{\formatdate{2}{10}{2013}}

%%First page's header%%%%%%%%%%%%%%%%%%%%%%
\fancypagestyle{firststyle} %%header settings for first pagr
{
	\fancyhead[L]{MTH \classNum{}(\classSec{}) \author{}\\
	\textbf{\examType{}} \date{}}
	\fancyhead[C]{\textbf{\nameType}}
	\setlength{\headsep}{14pt}
}
%%%%%%%%%%%%%%%%%%%%%%%%%%%%%%%%%%%%%%%%%%

%%Next page's header%%%%%%%%%%%%%%%%%%%%%%%
\setlist{nolistsep}
\lhead{MTH \classNum{}(\classSec{})} 
\chead{\textbf{\nameType}} 
\cfoot{\thepage{}~of~\pageref{LastPage}} 
\setlength{\headsep}{5pt} 
%%%%%%%%%%%%%%%%%%%%%%%%%%%%%%%%%%%%%%%%%%%

%%Various settings; see the preamble for other settings%%%%%
%%\setenumerate[0]{leftmargin=12pt,itemindent=0pt,labelwidth=0pt}
\renewcommand{\labelenumi}{\textbf{\arabic{enumi}.}}
\renewcommand{\labelenumii}{\textbf{\alph{enumii}.}}
\thispagestyle{firststyle}
%%%%%%%%%%%%%%%%%%%%%%%%%%%%%%%%%%%%%%%%%%
\newcommand{\xmax}{10}
\newcommand{\xmin}{-10}
\newcommand{\xstep}{0} %must be xmin+n
\newcommand{\ymax}{10}
\newcommand{\ymin}{-10}
\newcommand{\ystep}{0} %must be ymin+n

\newcommand{\Dspace}{\vspace{0cm}}

\textbf{Solve the radical equation, and check all proposed solutions}
\begin{enumerate}[resume]
\item $\sqrt{x+5}=7$
\begin{eqnarray*}
\left(\sqrt{x+5}\right)^{2} & = & 7^{2}\\
x+5 & = & 49\\
x & = & 44\checkmark
\end{eqnarray*}
\Dspace \vfill
\end{enumerate}
\textbf{Solve and check the equation.}
\begin{enumerate}[resume]
\item $x^{\nicefrac{3}{2}}=8$
\begin{eqnarray*}
\left(x^{\nicefrac{3}{2}}\right)^{\nicefrac{2}{3}} & = & 8^{\nicefrac{2}{3}}\\
x & = & \sqrt[3]{64}\\
x & = & 4\checkmark
\end{eqnarray*}
\Dspace \vfill
\item $(x+6)^{\nicefrac{5}{2}}-5=3$
\begin{eqnarray*}
(x+6)^{\nicefrac{5}{2}} & = & 8\\
\left((x+6)^{\nicefrac{5}{2}}\right)^{\nicefrac{2}{5}} & = & 8^{\nicefrac{2}{5}}=\left(2^{3}\right){}^{\nicefrac{2}{5}}=2^{\nicefrac{6}{5}}=2^{1}\cdot2^{\nicefrac{1}{5}}\\
x+6 & = & 2\sqrt[5]{2}\\
x & = & -6+2\sqrt[5]{2}\,\left(\approx-3.7026\right)\checkmark
\end{eqnarray*}
\Dspace \vfill
\end{enumerate}
\textbf{Solve the absolute value equation or indicate that the equation
has no solution.}
\begin{enumerate}[resume]
\item $\left|x+6\right|=7$
\begin{eqnarray*}
\left|x+6\right| & = & 7
\end{eqnarray*}
\[
\begin{array}{ccccccc}
x+6 & = & 7 & \mbox{or} & x+6 & = & -7\\
\Rightarrow x & = & 1\checkmark &  & \Rightarrow x & = & -13\checkmark
\end{array}
\]
\Dspace \vfill
\item $3\left|x+2\right|+4=2$
\begin{eqnarray*}
3\left|x+2\right| & = & -2\\
\left|x+2\right| & = & -\nicefrac{2}{3}\\
\left|x+2\right|\geq0 &  & \mbox{So there is NO solution}
\end{eqnarray*}
\Dspace \vfill
\end{enumerate}
\textbf{Solve the polynomial equation by factoring and then using
the zero product principle.}
\begin{enumerate}[resume]
\item $x^{3}+8x^{2}-x-8=0$
\begin{eqnarray*}
x^{2}(x+8)-1(x+8) & = & 0\\
(x^{2}-1)(x+8) & = & 0\\
(x+1)(x-1)(x+8) & = & 0\\
\Rightarrow x=-1,1,-8
\end{eqnarray*}
\Dspace \vfill
\end{enumerate}
\textbf{Determine if the given function is even, odd, or neither. }
\begin{enumerate}[resume]
\item \emph{$x^{4}+2x^{2}$. }

\begin{enumerate}[resume]
\item []It is even. Since the degree of each term is even. Also $f(-x)=(-x)^{4}+2(-x)^{2}=x^{4}+2x^{2}=f(x)$.\Dspace \vfill
\end{enumerate}
\end{enumerate}
\newpage\textbf{Give the domain and range of the relation, and determine
if it is a function. }{[}2 pts each{]}\vspace{-.25cm}

\begin{multicols}{2}
\begin{enumerate}[resume]
\item \vspace{0cm}$\left\{ (6,7),(10,-8),(9,6),(1,2)\right\} $

\begin{enumerate}[resume]
\item Domain:

\begin{enumerate}[resume]
\item []$\left\{ 6,10,9,1\right\} $
\end{enumerate}
\item Range:

\begin{enumerate}[resume]
\item []$\left\{ 7,-8,6,2\right\} $
\end{enumerate}
\item Function?

\begin{enumerate}[resume]
\item []Yes
\end{enumerate}
\end{enumerate}
\columnbreak
\item \vspace{0cm}$\left\{ (5,-4),(1,-3),(1,0),(10,-3),(26,-3)\right\} $

\begin{enumerate}
\item Domain:

\begin{enumerate}
\item []$\left\{ 5,1,10,26\right\} $
\end{enumerate}
\end{enumerate}
\begin{enumerate}[resume]
\item Range:

\begin{enumerate}[resume]
\item []$\left\{ -4,-3,0\right\} $
\end{enumerate}
\item Function?

\begin{enumerate}[resume]
\item []No, since $1$ is paired with both $-3$ and $0$.
\end{enumerate}
\end{enumerate}
\end{enumerate}
\end{multicols}

\textbf{Use the vertical line test to determine whether or not the
graph is a graph in which }\textbf{\emph{y }}\textbf{is a function
of }\textbf{\emph{x}}\textbf{. }{[}2 points each{]}
\begin{enumerate}[resume]
\item [\textbf{10.}]\setcounter{enumi}{10}\vspace{-.25cm}

\begin{multicols}{2}
\begin{enumerate}[resume]
\item \vspace{0cm}YES it is a function. A vertical line will pass through
the curve only once.
\end{enumerate}
\renewcommand{\xmax}{10}
\renewcommand{\xmin}{-10}
\renewcommand{\xstep}{0} %must be xmin+n
\renewcommand{\ymax}{10}
\renewcommand{\ymin}{-10}
\renewcommand{\ystep}{0} %must be ymin+n
%%%%%%%
\begin{tikzpicture} 
\begin{axis}[height=5cm, 
	width=5cm, 
	axis lines=middle,
	minor x tick num={4},
	minor y tick num={4},
	grid=both,
	%xtick={0,50,...,250},
	%ytick={0,10,20,...,40},
	%axis line style={-latex}, 
	xmin=\xmin,xmax=\xmax, 
	ymin=\ymin,ymax=\ymax,
	%restrict y to domain=-1:10,
	%no markers, 
	xlabel={$x$},ylabel={$y$}, 
	%every axis x label/.append style={anchor=north}, 
	%every axis y label/.append style={anchor=east}
	]
	
	%\addplot[color=black,solid,mark=*] coordinates {(0,1)} node[right] {\footnotesize{$(0,1)$}};

	\addplot[domain=\xmin:\xmax,thick,samples=100,draw=red]{1/3*(x)-2} node[pos=.525,right] {$$};
	\draw[very thick, blue, dashed] ({axis cs:2,0}|-{rel axis cs:0,0}) -- ({axis cs:2,0}|-{rel axis cs:0,1});
	
\end{axis}
\end{tikzpicture}
\begin{enumerate}[resume]
\item \vspace{0cm}No it is not a function by the vertical line test.
\end{enumerate}
\renewcommand{\xmax}{10}
\renewcommand{\xmin}{-10}
\renewcommand{\xstep}{0} %must be xmin+n
\renewcommand{\ymax}{10}
\renewcommand{\ymin}{-10}
\renewcommand{\ystep}{0} %must be ymin+n
%%%%%%%
\begin{tikzpicture} 
\begin{axis}[height=5cm, 
	width=5cm, 
	axis lines=middle,
	minor x tick num={4},
	minor y tick num={4},
	grid=both,
	%xtick={0,50,...,250},
	%ytick={0,10,20,...,40},
	%axis line style={-latex}, 
	xmin=\xmin,xmax=\xmax, 
	ymin=\ymin,ymax=\ymax,
	%restrict y to domain=-1:10,
	%no markers, 
	xlabel={$x$},ylabel={$y$}, 
	%every axis x label/.append style={anchor=north}, 
	%every axis y label/.append style={anchor=east}
	]
	
	%\addplot[color=black,solid,mark=*] coordinates {(0,1)} node[right] {\footnotesize{$(0,1)$}};
	\addplot[domain=\xmin:\xmax,red,thick,->] (1/2*(x-1)*(x-1)-3,x) node[pos=.525,right] {$$};
%%%Vertical Line test%%%%%%%%
	\draw[very thick, blue, dashed] ({axis cs:2,0}|-{rel axis cs:0,0}) -- ({axis cs:2,0}|-{rel axis cs:0,1});
%%%A big X%%%%%
	\addplot[thick,color=black] coordinates {(\xmin,\ymin) (\xmax,\ymax)};
	\addplot[very thick,color=black] coordinates {(\xmin,\ymax) (\xmax,\ymin)};
	
\end{axis}
\end{tikzpicture}

\end{multicols}
\end{enumerate}
\textbf{Use the provided graph to complete the problem.}
\begin{enumerate}[resume]
\item {[}3 points each{]}\vspace{-.3cm}

\begin{multicols}{2}
\begin{enumerate}
\item Determine the $x-$intercept(s) of $f(x)$.

\begin{enumerate}
\item []$\left\{ -5,2\right\} $
\end{enumerate}
\item Determine the $y-$intercept(s) of $f(x)$.

\begin{enumerate}
\item []$\left\{ 8\right\} $
\end{enumerate}
\item Determine the domain of $f(x)$.

\begin{enumerate}
\item []$\left(-\infty,\infty\right)$
\end{enumerate}
\item Determine the range of $f(x)$.

\begin{enumerate}
\item []$[-9,\infty)$
\end{enumerate}
\end{enumerate}
\columnbreak
\begin{enumerate}
\item []\renewcommand{\xmax}{10}
\renewcommand{\xmin}{-10}
\renewcommand{\xstep}{0} %must be xmin+n
\renewcommand{\ymax}{10}
\renewcommand{\ymin}{-10}
\renewcommand{\ystep}{0} %must be ymin+n
%%%%%%%
\begin{tikzpicture} 
\begin{axis}[height=5.5cm, 
	width=5.5cm, 
	axis lines=middle,
	minor x tick num={4},
	minor y tick num={4},
	grid=both,
	%xtick={0,50,...,250},
	%ytick={0,10,20,...,40},
	%axis line style={-latex}, 
	xmin=\xmin,xmax=\xmax, 
	ymin=\ymin,ymax=\ymax,
	%restrict y to domain=-1:10,
	%no markers, 
	xlabel={$x$},ylabel={$y$}, 
	%every axis x label/.append style={anchor=north}, 
	%every axis y label/.append style={anchor=east}
	]
	
	%\addplot[color=black,solid,mark=*] coordinates {(2,5)};
	\addplot[domain=\xmin:\xmax,thick,samples=100,draw=red,<->]{(x+4)*(x-2)} [yshift=0pt] node[pos=0.475,right] {$f(x)$};
		
\end{axis}
\end{tikzpicture}
\end{enumerate}
\end{multicols}
\end{enumerate}
\textbf{Evaluate the function at the given value of the independent
variable and simplify.}
\begin{enumerate}[resume]
\item $f(x)=-5x+1$; $f(-3)$
\begin{eqnarray*}
f(-3) & = & -5(-3)+1\\
 & = & 16
\end{eqnarray*}
\vfill \Dspace
\item $f(x)=x^{2}-x$; $f(x-2)$
\begin{eqnarray*}
f(x-2) & = & (x-2)^{2}-(x-2)\\
 & = & x^{2}-4x+4-x+2\\
 & = & x^{2}-5x+6
\end{eqnarray*}
\vfill \Dspace
\end{enumerate}
\textbf{Find the slope of the line that goes through the given points.}
\begin{enumerate}[resume]
\item $(2,-3)$ and $(-7,8)$.
\[
m=\frac{-3-(8)}{2-(-7)}=\frac{-11}{9}=-\frac{11}{9}
\]
\Dspace \vfill
\end{enumerate}
\textbf{Use the given conditions to write an equation for the line
in the indicated form.}
\begin{enumerate}[resume]
\item {[}8 points{]} Passing through $(5,8)$ with a slope of 3.

\begin{enumerate}[resume]
\item Point-slope form. $m=3$
\begin{eqnarray*}
y-8 & = & 3(x-5)
\end{eqnarray*}
\Dspace \vfill
\item Slope intercept form.
\begin{eqnarray*}
y-8 & = & 3(x-5)\\
y-8 & = & 3x-15\\
y & = & 3x-7
\end{eqnarray*}
\Dspace \vfill
\end{enumerate}
\item {[}8 points{]} Passing through $(3,8)$ and perpendicular to the line
$y=-\frac{1}{3}x+7$. 

\begin{enumerate}[resume]
\item Point-slope form. $m=3$ since $\frac{3}{1}$ is the negative reciprocal
of $-\frac{1}{3}$.
\begin{eqnarray*}
y-8 & = & 3(x-3)
\end{eqnarray*}
\Dspace \vfill
\item Slope intercept form.
\begin{eqnarray*}
y-8 & = & 3(x-3)\\
y-8 & = & 3x-9\\
y & = & 3x-1
\end{eqnarray*}
\Dspace \vfill\vspace{5cm}
\end{enumerate}
\end{enumerate}

\end{document}
